\documentclass[twoside,11pt]{article}

% Any additional packages needed should be included after jmlr2e.
% Note that jmlr2e.sty includes epsfig, amssymb, natbib and graphicx,
% and defines many common macros, such as 'proof' and 'example'.
%
% It also sets the bibliographystyle to plainnat; for more information on
% natbib citation styles, see the natbib documentation, a copy of which
% is archived at http://www.jmlr.org/format/natbib.pdf

\usepackage{jmlr2e}
\usepackage{amsmath}
\usepackage{amssymb}
\usepackage{bm}
\usepackage{preamble}

% todo
 \setlength {\marginparwidth }{2cm} % TODO: REMOVE THAT later
\usepackage{todonotes}
\definecolor{lightred}{RGB}{120,60,60}
\definecolor{darkred}{RGB}{255,0,0}
\newcommand{\proofread}[1]{{\color{lightred}#1}}
\newcommand{\rewrite}[2]{{\color{darkred}#1}}

\definecolor{gbcolor}{RGB}{220, 160, 180}
\newcommand{\gb}[2][]{\todo[color=gbcolor, caption={GB}, #1]{#2}}
\definecolor{jbcolor}{RGB}{180, 220, 200}
\newcommand{\jb}[2][]{\todo[color=jbcolor, caption={JB}, #1]{#2}}


% Definitions of handy macros can go here

\newcommand{\dataset}{{\cal D}}
\newcommand{\fracpartial}[2]{\frac{\partial #1}{\partial  #2}}

% Heading arguments are {volume}{year}{pages}{submitted}{published}{author-full-names}

\jmlrheading{1}{2000}{1-48}{4/00}{10/00}{Blanc}

% Short headings should be running head and authors last names

\ShortHeadings{The AI Trick for GDLVMs}{Blanc \& Bodelet}
\firstpageno{1}

\begin{document}

\title{Math summary for 
``The AI Trick: a Consistent Estimator for Deep Generalized Latent Variable Models''
}

\author{\name Guillaume Blanc \email guillaume.e.blanc@gmail.com \\
  \AND
  \name Julien Bodelet \email julien.bodelet@gmail.com \\
  \AND
  \name Benjamin Poilane \email 
}

\editor{Someone}

\maketitle


\section{Introduction}

The latent variable model is a \textit{generative model} $F_\bt$ for the pair $(\Y, \Z)$:
%
\begin{align}
    \begin{matrix}
        \hfill\Z &\sim& \pi(\z)\hfill\\
        \hfill\Y|\Z &\sim& f_\bt(\y|\z)\hfill,
    \end{matrix}
    \label{eqn:generative-model-general}
\end{align}
%
where $\Z$ is unobserved.
% 
The marginal density $f_\bt(\y)$ of $\Y$ is then obtained by integrating out the latent variable, 
%
\begin{equation}
    f_\bt(\y) = \int_{\mathcal Z} f_\bt(\y|\z)\pi(\z)d\z.
\end{equation}


Given $n$ independent data pairs $\{(\Y_i^0, \Z_i^0)\}_{i=1}^n$ from $F_{\bt_0}$, where $\Z_i^0$ are unobserved, an estimate of $\bt_0$ can, in principle, be obtained by maximizing the sample \textit{marginal log-likelihood}, 
%
\begin{equation}
\frac{1}{n}\sum_{i=1}^n l(\bt; \Y_{i}^0)= \frac{1}{n}\sum_{i=1}^n \log\left(\int_{\mathcal Z}f_\bt(\Y_{i}^0|\z)\pi(\z)d\z\right),
    \label{eqn:marginal-likelihood}
\end{equation}
The issue is that the integral is usually intractable.
We propose a framework that allows to overcome this issue.




\section{Motivation}

\subsection{The model}

A Generalized Deep Latent Variable Model (GDLVM)\gb[]{This doesn't exist, this is the model we want ( or not? ) to introduce. Perhaps we want to consider this model as a special case of our estimation strategy. We'll see that later.}, further specifies the conditional density $f_\bt(\y|\z)$, in two ways: (i) given the latent variable $\Z$, the elements $Y_{1}, \dots, Y_{p}$ of $\Y$ are independent, that is $Y_{j} \!\! \perp \!\!  Y_{j'}\vert \Z$, $\forall j\neq j'$. This allows to re-write $f_\bt(\y, \z) = \prod_{j=1}^p f_{\bt, j}(y, \z)$, where (ii) the conditional densities $f_{\bt, j}(\y|\z)$ belong to the exponential family of distributions, that is

\begin{equation}
   f_{\bt, j}(y|\z) = \exp(\{y\gamma - b_j(\gamma))\}/a_j(\phi_j) + c_j(y, \gamma)),
   \label{eqn:conditional-distribution}
\end{equation}
for some known functions $\{a_j, b_j, c_j\}_{j=1}^p$, arbitrary scale parameters $\{\phi_j\}_j^p$ and where $\gamma$ is called the cannonical parameter, which in turn entirely depends on $\bt$ and $\z$. 

\subsection{Profile likelihood}

In order to avoid the computation of the integral, we consider the extended likelihood:
\begin{equation}
    l(\bt, \z;\Y)= 
    \frac{1}{n}\sum_{i=1}^nl(\bt, \z_i;\Y_i)= 
    \frac{1}{n}\sum_{i=1}^n\log(f_\bt(\Y_{i}^0,\z_{i}))= 
    \frac{1}{n}\sum_{i=1}^n\log(f_\bt(\Y_{i}^0|\z_{i})\pi(\z_{i})).
    \label{eqn:extended-likelihood}
\end{equation} 
The (marginal) scores are defined by
$$
\bm\psi_n(\bt, \z) = \frac{1}{n}\sum_{i=1}^n
\bm\psi(\bt, \z_i, \y_i)
= \frac{1}{n}\sum_{i=1}^n \frac{\partial}{\partial\theta} l(\bt, \z_{i}; \Y_{i}). 
$$
In the case that  $\z$ is known, $\psi_n$ corresponds to the score of a GLM, and is Fisher consistent conditionally on $\z$:
$$
\E[\bm\psi_n(\bt, \z, \y)| \z] = 0,
$$
The profile likelihood treats $\z$ as a nuisance parameter. First, one obtains an estimator of $z$ given the parameter 
$\bt$ by solving the score equation:
$$
z_\bt: \bm\phi_n(\bt, \z) = 0.
$$
Then one solves the score by plugin the estimator of the nuisance parameter:
$$
\hat \bt_{\text{PL}}: \bm\psi_n(\bt, \z_\bt, \y) = 0.
$$
When it exists and $\bm\phi_n(\bt, \z) = \frac{1}{n}\sum_{i=1}^n \frac{\partial}{\partial\z} l(\bt, \z_{i}; \Y_{i})$,
the profile likelihood estimator is equivalent to the Maximum likelihood estimator.
However, while commonly used for mixed effect models, 
the profile likelihood fails for most latent variable models, since it is not Fisher consistent, i.e. $\E \bm\psi_n(\theta, \z_\theta, \y)\neq 0$.

\subsection{Partial profiling}

We propose a novel method which is based on partial profiling.
Consider the score equations for the exponential framework:
\begin{equation}
\bm\Psi_n(\bt, \z) = \sum_{j=1}^p\frac{\bm D_j(\z, \bt) (y_{ij}- \mu_j(\bt, \z_i))}{V_j(\z_i, \bt)\phi_j}.
\label{eqn:Psi}
\end{equation}
where, for all $j$, $\mu_j(\bt, \z) = \E[Y_{j}|\z] = g_j^{-1}(\eta_j(\bt, \z))$, $\bm D_j(\bt, \z) = \partial \mu_j(\bt, \z)/\partial \bt$, $V_j(\bt, \z)\phi_j = \text{Var}(Y_{j}|\z)$, $g_j$ is a link function, and $\phi_j$ is an additional variance parameter, and where the conditional expectation and variance are taken with respet to the conditional density $f_{\bt,j}(y|\z)$.
For simplicity we write
\begin{equation}
\bm\Psi_n(\bt, \z) = \frac{1}{n} \sum_{i=1}^n\sum_{j=1}^p \bm W_j(\z_i, \bt) (y_{ij}- \E_\bt(y_{ij}|z_i)).
\end{equation}
where we replaced $\mu_j(\bt, \z)$ by $\E_\bt(y|z)$.
We then do partial profiling, i.e. we replace the $\z_\bt$ only in the weights in the score:
$$
\tilde{\bm\Psi}_n(\bt, \z) = \frac{1}{n} \sum_{i=1}^n\sum_{j=1}^p \bm W_j(\z_{\bt,i}, \bt) (y_{ij}- \E_\bt(y_{ij}|\z_i)).
$$
In contrast to full profiling, the score function
$\tilde{\bm\Psi}_n(\bt, \z)$ still depend on $\bm z$ through the expectation.
However we argue that the dependence vanishes asymptotically.
Indeed, using the weak law of large number, we have 
\begin{align*}
\tilde{\bm\Psi}_n(\bt, \z) 
&=
\frac{1}{n} \sum_{i=1}^n\sum_{j=1}^p \bm W_j(\z_{\bt,i}, \bt) y_{ij} - 
\frac{1}{n} \sum_{i=1}^n\sum_{j=1}^p 
\bm W_j(\z_{\bt,i}, \bt) \E_\bt(y_{ij}|\z_i),\\    
&=
\frac{1}{n} \sum_{i=1}^n\sum_{j=1}^p \bm W_j(\z_{\bt,i}, \bt) y_{ij} - 
\E_\bt \sum_{j=1}^p \bm W_j(\z_{\bt,i}, \bt) y_{ij} + o_p(1)\\
&= \tilde{\bm\Psi}_n(\bt) +o_p(1).
\end{align*}
Basically, optimizing $\tilde{\bm\Psi}_n(\bt) $ is asymptotically equivalent to optimize $\tilde{\bm\Psi}_n(\bt, \z) $.
We thus use $\tilde{\bm\Psi}_n(\bt)$ as estimating equation for $\bt$.
We consider below the example of factor analysis which leads to explicite equations and a very fast algorithm.
Note that in general $E_\bt$ cannot be computed analytically.
We consider later the general scenario.


\subsection{Example with factor analysis}

Show off. Derive the partial profiled equations for FA.
Show numerical simulations and computing time.



\section{Algorithms for general cases}

\subsection{Algorithm with the map}

\subsection{Algorithm with the encoder}

Show algorithms with map and encoders.


\section{Asymptotic theory}

\subsection{Parametric case}

For consistency, we will need that the matrix $W_n(\y, \bt) \rightarrow W(\bt)$ and $W(\bt)$ is nonsingular (uniqueness).
Moreover, we that the derivatives are positive definites.
Basically, asymptotic theory is a particular case of Quasi-likelihood theory when the variance is misspecified.


\subsection{Nonparametric case}

\section{Simulations}

Simulations:
\begin{itemize}
    \item GLLVM parametric
    \item GLLVM nonparametric (Variational autoencoders)
    \item mixed models?
    \item Others? Bayesian?
\end{itemize}


\section{Application}

\subsection{Mnist}

\subsection{SNP}











\end{document}


